% !TXS template
\documentclass[french]{article}
\usepackage[T1]{fontenc}
\usepackage[utf8]{inputenc}
\usepackage{lmodern}
\usepackage[a4paper]{geometry}
\usepackage{babel}
\title{\textbf{EXERCICE DU CHAPITRE 2: Archéologie des Régimes de Vérité
		Numérique}}
\date{}
\begin{document}
\thispagestyle{empty}
\begin{center}
	\rule{1\textwidth}{2pt}\vspace{\baselineskip}
	\vspace{\baselineskip}\huge\textbf{THEORIES ET PRATIQUES DE L’INVESTIGATION  NUMERIQUE.}
	\rule{1\textwidth}{2pt}\vspace{\baselineskip}
	\vspace{\baselineskip}
	\vfill
	\textbf{Enseignant: M.MINKA THIERRY.}\\
	\vfill
	Par l'étudiante: CHEUTCHEULONGA HILARY PAULA\\
	Filière CIN 4\vfill
	\vspace{\baselineskip}
	\emph{\textbf{Année Academique: 2025-2026.}}
	\rule{1\textwidth}{2pt}\vspace{\baselineskip}
\end{center}
\maketitle

\section{Partie 1 : Analyse Historique et Epistémologique}
\subsection*{1.	Analyse Comparative des Régimes de Vérité}
a)	Période historique choisi est l’année 1990-2000 et 2010-2020\\
-	1990-2000 : R⃗ = (αT=0.2, αJ=0.4, αS=0.1, αP=0.3) = (0.2,0.4,0.1,0.3) ;\\
-	2010-2020 : R⃗ = (αT=0.5, αJ=0.2, αS=0.2, αP=0.1) = (0.5,0.2,0.2,0.1) ;\\
b)	Discontinuités foucaldiennes\\
-	Passage d’une vérité médiatique centralisée (journalisme, institutions) à une vérité distribuée (réseaux sociaux).\\
-	Changement du statut de l’expert, circulation des informations et multiplication des sources non vérifiées.\\
c)	Explication sociotechnique\\
-	Facteurs techniques : explosion des données, numérisation et accessibilité de l’internet.\\
-	Facteurs sociaux : confiance déplacée vers la donnée, développement des réseaux sociaux et du mobile.\\
-	Facteurs juridiques : adaptation tardive du droit.\\
d)	La transition a eu des signes dans les années 2000, mais le basculement post-2010 apparait révolutionnaire avec la viralité des réseaux sociaux et l’apparition de nouveaux régimes de véridiction.\\
\subsection*{2.	Etude de Cas Archéologique Foucaldienne}
a)	Choix de l’affaire\\
-	Affaire silk road (2011-2013)\\
b)	Analyse comme formation discursive au sens Foucault
-	Vérité algorithmique (blockchain, pseudonymes).\\
-	Discours de la liberté, de la décentralisation, de l’anonymat.\\
c)	Identification du « dicible » et le « pensable » à cette époque
-	Dicible : commerce anonyme et sécurisé sur internet, critique du monopole.\\
-	Pensable : internet comme espace de résistance, crypto-monnaie comme outils politiques.\\
d)	Comparaison avec une affaire contemporaine sous l’angle des régimes de vérité
-	Affaire actuelle : plateformes de crypto-monnaies (Binance, FTX)\\
-	Comparaison : passage d’un régime défensif (résistance à l’état) à un régime hybride (acceptation et institutionnalisation partielle), évolution du discours sur la régulation.\\

\section{Partie 2 : Modélisation Mathématique et Prospective}
\subsection*{4.	Vérification de l’accélération technologique}
-	1970-1990, 1990-2000, 2000-2010,2010-2020, 2020-Aujourdhui\\
-	Les intervalles entre régimes : 20 → 10 → 10 → 10 → 5 ans.\\
-	∆tn+1 = k · ∆tn \\
Avec n=1, ∆t2 = k · ∆t1 → 10=k*20 → K ≈ 0.5.\\
-	Prédiction : Si cette tendance se maintient, le prochain changement majeur pourrait survenir d'ici 5 à 7 ans vers 2030 (régime post-humain).\\
\subsection*{5.	Analyse du Trilemme CRO}
a)	Scores estimés :\\
-	1990–2000 : C=0.7, R=0.6, O=0.8\\
-	2000–2010 : C=0.8, R=0.9, O=0.7\\
-	2010–2020 : C=0.5, R=0.9, O=0.6\\
-	2020–... : C=0.9, R=0.7, O=0.4\\
b)	Compromis :\\
-	1990s → opposabilité.\\
-	2010s → fiabilité.\\
-	2020s → confidentialité.\\
c)	Futur : équilibre instable (Trilemme CRO non résolu).

\section{Partie 3 : Investigation Historique Appliquée}
\subsection*{6.	Reconstruction d’une investigation : Mitnick (1995)}
-	Époque : logs, surveillance manuelle, expertise humaine.\\
-	Aujourd’hui : IA de détection, corrélation automatisée.\\
-	Transition : du témoin expert à la preuve calculée : À l'époque, la vérité était construite principalement par des témoignages humains et des preuves matérielles ; aujourd'hui, elle repose davantage sur des données numériques et des algorithmes.\\
\subsection*{7.	Projet de recherche archéologique}
-	Trou : Un manque de documentation sur les débuts du P2P dans les années 2000, histoire non occidentale de la cyberforensique.\\
-	Hypothèse : Le développement des protocoles P2P a permis une nouvelle forme de vérité numérique basée sur le consensus collectif, pratiques hybrides locales.\\
-	Méthodologie : Analysez comment ces technologies ont redéfini le pouvoir et la vérité dans les communautés numériques, analyse RFC, archives, entretiens, méthode foucaldienne.\\
-	Analysez comment ces technologies ont redéfini le pouvoir et la vérité dans les communautés numériques.\\
\subsection*{8.	Analyse prospective des régimes futurs (2030–2050)}
-	Scénario 2030-2050 : Imaginons un monde où l'IA gère la plupart des décisions basées sur des données massives.\\
-	Régime de vérité : Ce régime pourrait être caractérisé par une transparence algorithmique et une confiance accrue dans les systèmes automatisés.\\
-	Conditions : Pour cela, il faudrait un cadre légal solide régulant l'usage de l'IA et garantissant la protection des données personnelles.\\
-	Méthodologie : neuro-forensique (analyse des signaux neuronaux).\\
-	Enjeux : éthique, libre arbitre, responsabilité algorithmique.\\

\end{document}
